\subsection{Shared intermediate sums for the bit network}
\label{sec:shared-sums}
We retain $h=46$, $v=\binom h3=15180$, $N=v^3$, $m=h^3=97336$, and
the rational form $I-J/9$ on $F=\mathbb Q^h$. Only the bit network changes;
the complex primitive is retained. Fix a common point $i$ of a neighboring
triple pair. On the other $n=h-1$ points, introduce inputs $x_{ab}$ indexed
by unordered pairs and compute all outputs
\[
 z_{cd}=\sum_{\{a,b\}\cap\{c,d\}=\varnothing}x_{ab}
 \qquad(\{c,d\}\in\tbinom{[n]}2).
\]
The sums are over $\mathbb F_2$. We use a circuit in which the two input
supports of every binary addition are disjoint. Equal formal sums are shared
as single nodes. Thus the circuit has no cancellation: a source contributes
to an output exactly when a directed path joins them, and that path is unique.

An explicit construction uses a balanced recursive split of the point set.
For a vector of disjoint-support values, compute the total and all sums
omitting one or two positions. The leave-two-out routine splits the vector
in half: exclusions in one half use that half's corresponding sum plus the
other half's total; split exclusions combine the two leave-one-out sums.
Totals use balanced binary addition. When only leave-one-out sums are needed,
use prefix and suffix sums, adding the prefix before a position to the suffix
after it. Zero terms are omitted throughout.

For pair inputs, return sums excluding zero, one, or two vertices. Recursively
compute these quantities for the left-left and right-right input pairs.
The cross pairs form a matrix. Its row sums and column sums feed the vector
routine just described. To omit one row and one column, apply leave-one-out
first within each row and then down each column of the resulting matrix.
For two excluded vertices in the left half, combine the recursive left-left
answer, the right-right total, and the cross-matrix row sum omitting those
two rows. The right-half case is symmetric. For one excluded vertex in each
half, combine the two recursive leave-one-out answers with the cross-matrix
answer omitting that row and column. Single exclusions and totals combine
the same three disjoint input regions. This proves correctness by induction.

Use the natural point order and splits at $\lfloor n/2\rfloor$, and identify
nodes whenever their formal input-support sets agree. Discard nodes not
feeding a requested leave-two-out output. The supplied implementation records
every support as an integer bit set and checks disjointness at every addition.
At $n=45$ the resulting circuit has
\[
 p=990\text{ inputs},\quad q=990\text{ outputs},\quad c=11566
 \text{ active binary additions}.
\]
Exact support comparison checks all 990 outputs, including all 893970 nonzero
matrix coefficients and all required zeros. The construction is a finite
explicit circuit; these counts make no optimality claim.

\subsection{Reversible embedding and the auxiliary-role count}
We next give a reusable transfer for such circuits. Treat each use of a node
by a later addition, or by a designated output, as an outgoing edge. Remove
unused nodes. At an input node, allocate one role per outgoing edge, select
one pivot as its input slot, and copy the pivot into the other outgoing roles.
At a binary addition node, the two incoming roles already exist. Reuse the
first as the pivot for one outgoing edge and allocate a fresh role for every
other outgoing edge. Add the second input into the pivot and then copy the
pivot into the other outgoing roles. The second incoming role is thereafter
inactive during the forward circuit. Each node operation is an invertible
linear map, with inverse obtained by undoing fanout and then the addition.

This rule identifies roles only along directed paths. In particular, two
simultaneous inputs never coincide. If there are $c$ addition nodes and $q$
designated outputs, the total number of outgoing edges is $2c+q$. One role
identification is made at each addition node, giving
\[
 R_{\rm local}=c+q=12556.
\]
Equivalently, an input contributes its outdegree in new roles and an addition
contributes its outdegree minus one. Inputs with one outgoing edge may be
regarded as identity gates, so every role participates in the finite mixer.
Let $L$ be the product of these node maps in topological order, $V$ the copy
into input slots, and $J$ the injection from designated output slots.
Then $JLV$ is precisely the disjoint-pair map. Each of the $h$ choices of
common point gets its own circuit, so
\[
 R_{\rm side}=h(c+q)=577576
\]
side roles suffice per invocation.

Arbitrary initial scratch is handled by the same transparent schedule as for
the rectangle circuit. With $G,R$ denoting the old central gather and scatter,
the chronological operations are
\[
 L,\ J,\ L^{-1},\ R,\ V,\ G,\ R,\ L,\ J,\ L^{-1},\ G,\ V.
\]
An initial side vector $z$ contributes $JLz$ in the early injection and
$JL(z+Vx)$ in the later one, leaving exactly $JLVx$. The final inverse mixer
and copy restore $z$; centers are restored as before. The complete side map
is the intersection-one neighbor matrix, since each neighbor has a unique
common point. Thus $JLV+RG=I$ and the invocation is $y\gets y+x$ on all inputs.

\subsection{Frames from source and target supports}
Write the original tensor decomposition at a stage as
$A=B\mathbin\perp P$, with $P$ and the future factor $Q$ nondegenerate lines.
Suppress $Q$ and put
\[
 D_0=B\otimes F,\qquad D_1=A\otimes F,\qquad
 D_U=(B\otimes F)\mathbin\perp(P\otimes U).
\]
During the middle forward application of $L$, assign each node the span $U$
of the physical source triple indicators reaching it. At an input node this
is its source line. Along every dataflow edge the source-support sets grow,
so these spans are nested. Along the physical roles, which follow paths
through the circuit, the labels are therefore increasing. A role allocated
at a branching node enters from $D_0$; a retired input role subsequently
grows to $D_1$ in the final cleanup.

Every triple in a source-support span contains its circuit's common point
$i$. For any vector $u$ in such a span, $\sum_j u_j=3u_i$, and hence
\[
 \langle u,u\rangle=\sum_{j\ne i}u_j^2>0\qquad(u\ne0).
\]
Thus all these spans are nondegenerate. At an output, every contributing
source triple is a neighbor of the target triple $T$, giving
$U\subset t_T^\perp$. The output role can consequently grow to the original
physical target frame without decreasing rank. Every node gate has one
common frame on all its incoming and outgoing roles, as required by the
common-frame identity.

The reversed stage requires a different span assignment. Reverse and invert
the scalar schedule with logical source $Y$ and target $X$. In the middle
$L^{-1}$, label a node by the span of the physical $X$ target indicators
reachable from it in the forward dataflow graph. These descendant sets grow
along reversed edges. Every such triple still contains the same point $i$,
so its span is nondegenerate by the same norm identity. At an original input,
all its reachable targets are neighbors of its physical $Y$ triple; the span
therefore lies in that triple's orthogonal complement. This verifies the
reverse middle computation as well as its endpoints. The source-support
frames of the forward computation are not incorrectly reused in reverse.

All early mixers and their inverses use $D_0\otimes Q$, and all late cleanup
mixers and inverses use $D_1\otimes Q$. The data-copy and injection gates are
grouped by data triple, with the same individual line and orthogonal-complement
frames as in the rectangle construction. Physical $X$ grows from
$A\otimes\langle t\rangle\otimes Q$ to $(A\otimes F)\otimes Q$;
physical $Y$ grows from $B\otimes\langle t\rangle\otimes Q$ to
$((B\otimes F)\mathbin\perp(P\otimes t^\perp))\otimes Q$.
These are exactly the original stage boundaries. The only decreasing edges
remain the $h$ central returns per invocation, each losing dimension $h$.
Consequently $L_{\rm dec}=3v^2h^2$ is unchanged.

\subsection{Stage sharing and rank accounting}
Use a bijection $\pi$ of triples with $|T\cap\pi(T)|=1$. One explicit
choice pairs the ground points consecutively. For a triple with a full pair
and a singleton, keep the singleton and cycle the full pair among all pairs
other than the singleton's pair. Otherwise keep the point in the least
indexed pair and flip the other two to their partners. This is a bijection
on the two disjoint classes, with the stated intersection property.

Share the entire auxiliary bank of the stage-1 invocation with fixed triples
$(A,B)$ with that of stage 3 at $(B,\pi(A))$, keeping each local role index.
Every invocation restores arbitrary scratch, and the sharing is bijective.
The joining labels are
\[
 E=F\otimes\langle t_A\rangle\otimes\langle t_B\rangle,
 \qquad H=\langle t_B\otimes t_{\pi(A)}\rangle^\perp\otimes F.
\]
They are nondegenerate and $E\subset H$ since $t_A\perp t_{\pi(A)}$.
Replacing the two former terminal edges by this join saves exactly $m$ in
edge rank per removed role. Thus
\begin{align*}
 W_{\rm b}^{\rm dag}&=2N+2v^2(R_{\rm side}+h)=273201575169600,\\
 s_{\rm b}^{\rm dag}&=W_{\rm b}^{\rm dag}m-N+6v^2h^2
                         =26592347948314104000,\\
 W_{\rm b}^{\rm dag}m-s_{\rm b}^{\rm dag}&=572394081600>0.
\end{align*}

\begin{proposition}[Shared-computation bit interface]\label{prop:dag-bit-interface}
The new bit network exchanges the two scalar data banks, restores every
auxiliary value, and satisfies the original all-role endpoint identity with
$W_{\rm b}^{\rm dag}$ roles and total edge rank $s_{\rm b}^{\rm dag}$.
\end{proposition}
\begin{proof}
The transparent schedule proves the scalar shear and arbitrary restoration.
Forward, inverse, forward invocations give the bank exchange. The support
labels and their reverse counterparts establish all common frames and nested
side transitions. Central decreases are unchanged. Telescoping therefore gives
$W_{\rm b}^{\rm dag}m-2N+2L_{\rm dec}$ absolute label change; the original
rational source correction adds $N$. Every surviving auxiliary role still
has endpoint matrices zero and $I_m$ and is fixed by the scalar permutation;
the data endpoints are unchanged. The original common-frame and rational-shear
transfer consequently applies. All compiled gates and bases are fixed finite
data, independent of the eventual input length.
\end{proof}

\subsection{Explicit rational bounds for the two recurrences}
Put
\[
 a_{\rm b}=\frac{187}{10^{11}},\qquad
 a_{\rm c}=\frac9{500000000000},\qquad L_0=\frac{5743}{500}.
\]
The exact deficits satisfy
\[
 \eta_{\rm b}^{\rm dag}=\frac{27}{1254369032}>a_{\rm b}L_0,
 \qquad \eta_{\rm c}=\frac7{22253827054}>a_{\rm c}L_0.
\]
The retained rational logarithm enclosure proves $\log m<L_0$. Therefore
$m^{1-a}>m(1-aL_0)>m(1-\eta)$ for either corresponding pair, and we may take
\begin{equation}\label{eq:explicit-motif-exponents}
 \tau=1-a_{\rm b},\qquad\sigma=1-a_{\rm c}.
\end{equation}
This gives $s_{\rm b}^{\rm dag}/W_{\rm b}^{\rm dag}<m^\tau$ and the
unchanged complex bound $s_{\rm c}/W_{\rm c}<m^\sigma$.
